\chapter{Build, Run, and Reproducibility Record}
\label{app:reproducibility}

This appendix records how to obtain, build, run, and reproduce the system and
the evidence reported in the thesis; the full user-facing instructions are in
Chapter~\ref{ch:user-documentation}.

\section{Repository, revision, environment, and install}

The source is in a private GitHub repository to which the supervisor has
access, and the complete source is part of the submission archive:

\url{https://github.com/thonmay/ai-dynamic-graph-algorithms}

The thesis describes the code at commit \texttt{7895107}. The two later commits
that touch the code add only the point-of-use AI-attribution comments
(\texttt{e39d4ab}; Python ASTs identical, UI bundle byte-identical) and fix one
flaky test (\texttt{e527d08}); all other later commits touch only the thesis
documents, the evidence logs and measurement scripts. The same-revision execution
record was run on \texttt{e527d08}, and the submission manifest
(\texttt{SUBMISSION\_MANIFEST.md}, Appendix~\ref{app:declarations}) will tie the
final revision to the PDF and the archive by hash. The supported environment is
Linux x86-64 with unprivileged user namespaces (the sandbox jail requires them),
Python 3.11 or later, and Node 24 only to rebuild the user interface; the thesis
PDF is built with TeX Live 2026. No database is used. The editable install makes
\texttt{evoalgo} a thin wrapper around \texttt{src/evoalgo/cli.py}:

\begin{lstlisting}[language=bash, basicstyle=\ttfamily\small]
git clone https://github.com/thonmay/ai-dynamic-graph-algorithms
cd ai-dynamic-graph-algorithms
python3 -m venv .venv
.venv/bin/pip install -e '.[api,ui]'
# rebuild the web bundle only if the UI source changed (Node 24 on the path):
cd ui && npm ci && npx vitest run && npx tsc --noEmit -p . && npm run build
\end{lstlisting}

\section{Run the test suites, self-tests, and campaigns}

The fast Python suite excludes the 5 slow corpus tests (about 20 minutes). The
self-test drives the real verifier over the full bound and prints
\texttt{ALL PASS} on success; the second form adds two relabelled copies of every
\texttt{nauty} class, the universe of the final grade (Section~\ref{sec:test-selftest}).
A campaign is a grid of (mode, problem, seed) cells defined by a YAML plan;
\texttt{run} is resumable and \texttt{report} reads only the run folders, so it
can be re-run at any time (see the header comments of \texttt{scripts/campaign.py}).
The cross-check, the random-update check and the chunk measurement are the
helper scripts named in Sections~\ref{sec:test-selftest} and
\ref{sec:test-reliability}; \texttt{execution\_record.sh} reproduces the whole
record of Section~\ref{sec:test-strategy} from a fresh clone. The figures are
standalone TikZ documents built before the thesis:

\begin{lstlisting}[language=bash, basicstyle=\ttfamily\small]
.venv/bin/pytest -q -m "not slow" --basetemp=<scratch>/pt -p no:cacheprovider
.venv/bin/pytest -q -m slow --basetemp=<scratch>/pt -p no:cacheprovider
cd ui && npx vitest run && npx tsc --noEmit -p .
.venv/bin/evoalgo verify --problem dynamic_bfs --n-max 5
.venv/bin/evoalgo verify --problem dynamic_dfs --n-max 5
.venv/bin/evoalgo verify --problem dynamic_dfs --n-max 5 --relabels 2
.venv/bin/evoalgo verify --problem dynamic_lbfs --n-max 5
.venv/bin/python scripts/campaign.py run campaign/eval1.yaml --parallel 4
.venv/bin/python scripts/campaign.py report campaign/eval1.yaml
.venv/bin/python scripts/dfs_certificate_crosscheck.py
.venv/bin/python scripts/chunk_memory.py dynamic_dfs dynamic_bfs
scripts/execution_record.sh <repo-url> <revision> <out-dir> python3
cd thesis-doc && make -C figures && latexmk -pdf main.tex
\end{lstlisting}

\section{Evidence artefacts, dependencies, and what is not reproducible}

All campaign evidence is committed under \texttt{campaign/}: the evaluation
results are \texttt{eval1/}, \texttt{eval2/}, and \texttt{eval3/}; the self-test
logs are \texttt{selftest\_bfs\_n5.log}, \texttt{selftest\_dfs\_n5.log},
\texttt{selftest\_dfs\_n5\_relabels2.log} and \texttt{selftest\_lbfs\_n5.log};
the checks beyond the self-test are \texttt{dfs\_certificate\_crosscheck.log},
\texttt{yang\_random\_check.log} and \texttt{chunk\_memory.log}; the
same-revision execution record is \texttt{execution\_record/}; the
user interface screenshots are \texttt{ui\_check/}; and the plans are
\texttt{eval1.yaml}, \texttt{eval2.yaml}, and \texttt{eval3.yaml}.

New runs require an OpenAI-compatible LLM endpoint; the mock provider works
offline and is used by the test suites, and recorded runs replay offline in the
user interface. The verifier, the sandbox, and the parent selection are
deterministic and reproducible. What is not reproducible bit for bit is the LLM
sampling: the model's sampling is not seeded, so a rerun of a campaign cell
produces different candidates and timings, even though the verifier's verdict on
any given candidate is reproducible.
